\section{Reference implementation}
\label{app:impl}

The listing below implements the algorithms of
\Cref{alg:rotate,alg:rebalance,alg:insert,alg:delete} directly, so that each
procedure in the text has one counterpart here. Heights are stored at the nodes
and recomputed by \texttt{update} after any change, following the convention of
\Cref{sec:prelim}: a leaf has height $0$ and the empty tree has height $-1$,
which is what lets \texttt{height} return $-1$ for \texttt{None} without a
special case anywhere else.

Both update procedures are written recursively and return the new root of the
subtree they were given. The caller relinks that return value, which is how a
rotation that changes the root of a subtree is propagated to its parent without
storing parent pointers.

\lstset{
  language        = Python,
  basicstyle      = \ttfamily\footnotesize,
  keywordstyle    = \bfseries,
  commentstyle    = \itshape\color{rmid},
  stringstyle     = \color{rmid},
  numbers         = left,
  numberstyle     = \tiny\color{rfaint},
  numbersep       = 8pt,
  frame           = leftline,
  rulecolor       = \color{rfaint},
  framesep        = 8pt,
  xleftmargin     = 18pt,
  showstringspaces = false,
  breaklines      = true,
  columns         = fullflexible,
}

\begin{lstlisting}
class Node:
    __slots__ = ("key", "left", "right", "h")

    def __init__(self, key):
        self.key = key
        self.left = None
        self.right = None
        self.h = 0            # a leaf has height 0


def height(n):
    return -1 if n is None else n.h      # the empty tree has height -1


def update(n):
    n.h = 1 + max(height(n.left), height(n.right))


def bf(n):
    """Balance factor: positive when the left subtree is the taller one."""
    return height(n.left) - height(n.right)


# --------------------------------------------------------------- rotations

def rotate_right(z):
    y = z.left
    z.left = y.right          # the middle subtree changes parent
    y.right = z
    update(z)                 # z is now the lower of the two
    update(y)
    return y                  # y is the new root of this subtree


def rotate_left(z):
    y = z.right
    z.right = y.left
    y.left = z
    update(z)
    update(y)
    return y


def rebalance(n):
    """Restore the invariant at n, returning the node that now roots it."""
    update(n)
    if bf(n) > 1:                         # left side too tall
        if bf(n.left) < 0:                # LR: straighten to LL first
            n.left = rotate_left(n.left)
        return rotate_right(n)
    if bf(n) < -1:                        # right side too tall
        if bf(n.right) > 0:               # RL: straighten to RR first
            n.right = rotate_right(n.right)
        return rotate_left(n)
    return n                              # already legal


# ----------------------------------------------------------------- updates

def insert(n, key):
    if n is None:
        return Node(key)
    if key < n.key:
        n.left = insert(n.left, key)
    elif key > n.key:
        n.right = insert(n.right, key)
    else:
        return n                          # the key is already present
    return rebalance(n)


def delete(n, key):
    if n is None:
        return None
    if key < n.key:
        n.left = delete(n.left, key)
    elif key > n.key:
        n.right = delete(n.right, key)
    else:
        if n.left is None:
            return n.right
        if n.right is None:
            return n.left
        s = n.right                       # in-order successor
        while s.left:
            s = s.left
        n.key = s.key
        n.right = delete(n.right, s.key)
    return rebalance(n)


def search(n, key):
    while n is not None and n.key != key:
        n = n.left if key < n.key else n.right
    return n
\end{lstlisting}

One detail is worth drawing out, because it is where the difference between
\Cref{thm:insertone} and \Cref{thm:deletebound} appears in the code. The two
update procedures look alike: both recurse down, both call \texttt{rebalance} on
the way out at every node along the path. The difference is not in the control
flow but in what \texttt{rebalance} does when it gets there. After an insertion,
the first rotation restores the original subtree height, so every call above it
finds a balance factor already in range and returns immediately. After a
deletion, a rotation can leave the subtree one level shorter, so the next call up
may find work to do. The same loop produces one repair in the first case and up
to $h$ in the second.
